% !TeX program = lualatex
\documentclass[12pt]{article}
\usepackage{fontawesome5}
\usepackage{multicol}
\usepackage[margin=1.5cm]{geometry}
\usepackage[hidelinks]{hyperref}

\begin{document}

\section*{Font Awesome piktogramų katalogas}

\noindent Naudojimas: \verb|\faIcon{pavadinimas}|. Piktogramos pavadinimas
pateiktas šalia simbolio.\medskip

\small
\setlength{\parindent}{0pt}
\setlength{\columnsep}{1cm}

\directlua{
  local file = kpse.find_file("fontawesome5-mapping.def")
  if not file then
    error("Nepavyko rasti Font Awesome piktogramu saraso.")
  end
  local input = assert(io.open(file, "r"))
  local icons = { free = {}, brands = {} }
  local left = string.char(123)
  local right = string.char(125)
  for line in input:lines() do
    if line:find("__fontawesome_def_icon:nnnnn", 1, true) then
      local fields = {}
      local start = 1
      while true do
        local a, b = line:find(left, start, true)
        if not a then break end
        local c, d = line:find(right, b + 1, true)
        if not c then break end
        table.insert(fields, line:sub(b + 1, c - 1))
        start = d + 1
      end
      local name = fields[2]
      local family = fields[3]
      if name and family then
        local group = family:sub(1, 4) == "free" and "free"
          or family:sub(1, 6) == "brands" and "brands"
        if group then
          table.insert(icons[group], name)
        end
      end
    end
  end
  input:close()
  table.sort(icons.free)
  table.sort(icons.brands)
  local slash = string.char(92)
  local function print_group(title, list)
    tex.sprint(slash .. "section*{" .. title .. "}")
    tex.sprint(slash .. "begin{multicols}{3}")
    for _, name in ipairs(list) do
      if name == "link" then
        tex.sprint(slash .. "href{https://drive.google.com/file/d/1K_VWxzg5jwKyv5-phhwlEzPQp0TV2cES/view?usp=sharing}{"
          .. slash .. "faIcon{link}} "
          .. slash .. "texttt{link}"
          .. slash .. slash)
      else
        tex.sprint(slash .. "faIcon{" .. name .. "} "
          .. slash .. "texttt{" .. name .. "}"
          .. slash .. slash)
      end
    end
    tex.sprint(slash .. "end{multicols}")
  end
  print_group("Paprastos piktogramos", icons.free)
  print_group("Prekių ženklų piktogramos", icons.brands)
}

\end{document}